\scanpage{40}
\mspage{36}
% archive sheet number: 39
% source lines 1--7; continuation from scan 39
alors qu'un objet \inserted{$X$} de $M$ est un \deleted{\unsure{obj}} élément (de nature particulière) de $\unsure{\mathcal P|\mathscr P}(\mathcal P(\unsure{L|\mathcal L}))$, une multistrate $\widetilde X$ de $M$ doit être \unsure{considéré|considérée}, en bonne logique, comme un élément de $\mathcal P(\mathcal P(\mathcal P(L))$ --- un ensemble d'ensembles de parties de $L$.

% source lines 8--23
De plus en plus, $M$ apparaissant comme l'ensemble de base privilégié, les \uline{figures} $F$ apparaissent comme des parties de $M$. Dans ce point de vue, il n'y a donc pas lieu d'introduire un $\widetilde F$, qui ici n'est autre que $F$. Donc il n'y a pas lieu de parler du ``déploiement'' d'une figure. C'est à l'usage qu'on verra s'il y a lieu, par la suite, d'autonomiser $\unsure{\mathcal F|\mathscr F|\mathfrak F}$ \unsure{p.v.|p.r.} à $M$, par l'introduction d'une structure mixte sur $(M,\mathcal F)$, avec une relation d'incidence $\trianglelefteq$ entre $M$ et $\mathcal F$ --- auquel cas la notation $\widetilde F$ reprend un sens et ses droits.
